%=====================================================================
\section{QCD Phenomenology from $\mathbb{C}^5$}
\label{app:qcd}
%=====================================================================

This appendix develops the non-perturbative aspects of the strong force from the spatial sector $\mathbb{C}^3 \subset \mathbb{C}^5$. The key results --- confinement, asymptotic freedom, the deconfinement transition, and the strongly-coupled quark-gluon plasma (sQGP) --- all follow from the rank structure of the spatial Gram sub-matrix.

\subsection{The moment map: $\mathbb{CP}^4 \to \Delta^4$}

Each normalized state $\psi \in \mathbb{CP}^4$ maps to its probability vector:
\begin{equation}
\mu(\psi) = (|\psi_0|^2, \ldots, |\psi_4|^2) \in \Delta^4
\end{equation}
where $\Delta^4$ is the standard 4-simplex ($\sum_a |\psi_a|^2 = 1$). This is the moment map for the $\text{U}(1)^5$ torus action on $\mathbb{CP}^4$. The $(3,2)$ split gives:
\begin{equation}
\mu_S = \sum_{a=0}^{2}|\psi_a|^2 \quad (\text{spatial weight}), \qquad \mu_T = \sum_{a=3}^{4}|\psi_a|^2 \quad (\text{temporal weight}), \qquad \mu_S + \mu_T = 1.
\end{equation}

\subsection{Confinement as $\mathbb{C}^3$ saturation}

The spatial overlap between vertices $i$ and $j$:
\begin{equation}
G^S_{ij} = \sum_{a=0}^{2} \psi^*_{i,a}\,\psi_{j,a}
\end{equation}
has $|G^S_{ij}| \leq \mu_S(i)^{1/2}\,\mu_S(j)^{1/2}$ by Cauchy--Schwarz. The spatial Gram matrix $G^S$ has rank at most $n_S = 3$.

\textbf{Confinement criterion:} When 3 vertices $\{a, b, c\}$ span all of $\mathbb{C}^3$ (i.e., $\text{rank}(G^S_{abc}) = 3$), the spatial sector is \emph{saturated}. No additional vertex can be spatially independent. Any fourth vertex must be a linear combination of $\{a, b, c\}$ in the spatial sector --- it is ``confined'' to the color space already defined.

This is the lattice analogue of color confinement: quarks cannot exist as free particles because the $\mathbb{C}^3$ sector is always locally saturated. The confinement scale is not a free parameter but a geometric property of $n_S = 3$.

\subsection{Asymptotic freedom as $\mathbb{C}^5$ dilution}

At high energies (short distances), the full $\mathbb{C}^5$ becomes relevant. The effective coupling strength is:
\begin{equation}
\alpha_\text{eff}(Q) \propto \frac{n_S}{d} = \frac{3}{5}
\end{equation}
at the GUT scale, where all 5 dimensions participate equally. As energy decreases, the $(3,2)$ split freezes and the strong force sees only $\mathbb{C}^3$:
\begin{equation}
\alpha_\text{eff}(Q) \propto \frac{n_S}{n_S} = 1 \qquad (\text{low energy: confined}).
\end{equation}
The coupling grows from $3/5$ to $1$ as energy decreases --- this is asymptotic freedom.

In DRLT, asymptotic freedom is equivalent to Pachner coarse-graining: at lower resolution (fewer effective vertices), the spatial sector occupies a larger fraction of the available Hilbert space, increasing the effective coupling. The standard QCD $\beta$-function $b_3 = -7$ is the derivative of this geometric transition.

\subsection{Deconfinement and the sQGP}

\subsubsection{The deconfinement transition}

At temperature $T > T_c$ (energy above the confinement scale), the thermal fluctuations overwhelm the spatial rank constraint. The spatial overlap $G^S_{ij}$ becomes large for many vertex pairs simultaneously, and $\mathbb{C}^3$ is globally populated rather than locally saturated.

The order parameter is the Polyakov loop analogue:
\begin{equation}
L = \frac{1}{N}\sum_i \mu_S(i)
\end{equation}
which transitions from $\langle L\rangle \approx 3/5$ (confined, spatial weight $= n_S/d$) to $\langle L\rangle \to 1$ (deconfined, spatial sector dominates). Numerical simulations show this transition occurs near $N_\text{eff} \approx 21 = 3 \times 7$ ($= n_S \times |b_3|$), consistent with the QCD crossover temperature.

\subsubsection{Strongly-coupled QGP}

Above $T_c$ but below the GUT scale, the quark-gluon plasma is deconfined but $W_{ij}$ remains large (strong correlations persist). This is the sQGP regime:
\begin{itemize}
\item Deconfined: $\text{rank}(G^S) = 3$ is no longer locally saturated.
\item Strongly coupled: $\langle W\rangle \gg 1/d$ (vertices are not random).
\item Nearly perfect fluid: the ratio $\eta/s$ approaches the KSS bound $1/(4\pi)$.
\end{itemize}

The KSS bound has a geometric interpretation in DRLT: the minimum viscosity corresponds to the maximum entropy density, which occurs when the Wishart eigenvalues are maximally spread --- a configuration with $\eta/s = 1/(4\pi)$ is the most ``democratic'' distribution of correlations across the lattice.

\subsection{The QCD phase diagram from $\mathbb{C}^3$}

\begin{center}
\begin{tabular}{lccl}
\hline
Regime & $G^S$ rank & $\langle W_S\rangle$ & Physics \\
\hline
Confined (hadrons) & 3 (locally saturated) & High (local) & Color singlets only \\
Crossover ($T \sim T_c$) & 3 (globally spreading) & Moderate & Deconfinement \\
sQGP ($T > T_c$) & 3 (global) & Still high & Nearly perfect fluid \\
Asymptotically free & 3 (diluted in $\mathbb{C}^5$) & Low ($\to 3/5d$) & Weakly coupled \\
\hline
\end{tabular}
\end{center}

The entire QCD phase structure follows from one integer: $n_S = 3$.
